\begin{helper}
Fix one support-label pair \(B,J\) in a history cell, and write \(U\) for
the terminal coordinate set with nodup reveal order \(o\).  Let
\(\mathcal P\) be the transcript family.  For each realized
branch/transcript pair \((\beta,\tau)\), write
\[
P_\tau,\qquad A_\tau,\qquad Z_\tau,\qquad F_{\beta,\tau}
\]
for the forced-present set, forced-absent set, selected absent set, and the
fixed present support.  Assume the realized geometry
\[
F_{\beta,\tau}\subseteq P_\tau,\qquad
P_\tau,A_\tau\subseteq U,\qquad
P_\tau\cap A_\tau=\emptyset,\qquad
Z_\tau\subseteq A_\tau ,
\]
the forward compatibility consequences, and functional disjointness for
complete assignments compatible with two realized cells.

Assume in addition the two fixed-\((B,J)\) residual compatibility inputs
supplied by the branch/transcript scan construction.  First, every residual
assignment in the cylinder
\[
P_\tau\setminus F_{\beta,\tau}\subseteq A_{\rm res},\qquad
A_{\rm res}\cap(A_\tau\setminus Z_\tau)=\emptyset
\]
reconstructs to a compatible complete assignment
\[
(A_{\rm res}\setminus Z_\tau)\cup F_{\beta,\tau}.
\]
Second, the same residual assignment cannot reconstruct compatibly to two
different realized pairs.  Then there are two option-valued functions: a
residual scan on residual assignments and a reconstructed complete-assignment
scan.  The reconstructed scan output is equivalent to complete compatibility,
the residual scan output is equivalent to compatibility after deleting
selected absences and reinserting fixed support, the two scans agree on each
residual cylinder, the reconstructed scan returns the cylinder owner there,
and the residual scan is functional on a fixed residual assignment.
\end{helper}
\begin{proof}
The point of the statement is that scan data are constructed only after the
missing compatibility inputs are present.  The residual-completeness and
single-owner inputs are the fixed-\((B,J)\) consequences of the
blue-first/red-second reveal-order construction isolated in
\(\noderef{FixedSetHistoryCellRedBranchTranscriptCells}\), with the
reconstruction stability coming from
\(\noderef{FixedSetHistoryCellRISIResidualScanStability}\) and the
first-mismatch uniqueness from
\(\noderef{FixedSetHistoryCellFixedBJResidualOverlapConsistency}\).

Define the reconstructed complete-assignment scan as follows.  On a complete
terminal assignment \(A\), if there is a realized pair
\((\beta,\tau)\) such that \(A\) is compatible with \((\beta,\tau)\), choose
one such pair and output it; otherwise output \(\bot\).  Complete-assignment
functional disjointness makes this choice unique, so whenever compatibility
with a realized \((\beta,\tau)\) holds, the reconstructed scan outputs
\((\beta,\tau)\), and every reconstructed-scan output gives the corresponding
compatibility certificate.

Define the residual scan similarly, but using reconstructed compatibility:
on \(A_{\rm res}\), choose the unique realized pair \((\beta,\tau)\) for
which
\[
(A_{\rm res}\setminus Z_\tau)\cup F_{\beta,\tau}
\]
is compatible with \((\beta,\tau)\), if such a pair exists.  The residual
single-owner input says that this choice is unique.  Thus residual-scan
output implies compatibility of the reconstructed complete assignment, and
reconstructed compatibility implies the corresponding residual-scan output.

Now suppose \(A_{\rm res}\) lies in the residual cylinder of a realized
\((\beta,\tau)\).  The residual-completeness input gives compatibility of
\((A_{\rm res}\setminus Z_\tau)\cup F_{\beta,\tau}\) with
\((\beta,\tau)\).  By the two definitions just made, both the residual scan
on \(A_{\rm res}\) and the reconstructed scan on that reconstructed complete
assignment output \((\beta,\tau)\).  Hence the scans agree on the cylinder,
and the reconstructed scan returns the cylinder owner.  Finally, if the
residual scan has two outputs on the same residual assignment, they are both
the value of the same option-valued function, so the branch labels and
transcripts are equal.  These are exactly the advertised fixed scan-interface
clauses.
\end{proof}
